\section{Proof of Impossibility and Approximate Impossibility}
\label{app:impossiblility}

In this section, we prove that it is impossible to satisfy multiple
equal-cost constraints while simultaneously satisfying calibration.
We will first prove this in an exact sense, and then show that
the result holds approximately as well.

\subsection{Exact Impossibility Theorem}

\newtheorem*{thm:impossibility}{Theorem \ref{thm:impossibility} (Restated)}
\begin{thm:impossibility}
  Let $h_1$ and $h_2$ be calibrated classifiers for $G_1$ and $G_2$ with equal
  cost with respect to
  $\fairconstname$.
  If $\mu_1 \ne \mu_2$, and if $h_1$ and $h_2$ have equal cost with respect to
  $\fairconstname'$, then $h_1$ and $h_2$ must be perfect classifiers.
\end{thm:impossibility}
%
\begin{proof}
  First, observe that the perfect classifier always satisfies any equal-cost
  constraint simply because if $\fp{t} = \fn{t} = 0$,
  $\fairconst{h_t}{t} = 0$. Moreover, the perfect classifier is
  always calibrated.

  For any classifier, as shown in \cite{kleinberg2016inherent}, \fp{h_t} and
  \fn{h_t} are linearly related by \eqref{eqn:linear-calib-sup}.
  Furthermore, each equal-cost constraint is linear in \fp{h_t} and
  \fn{h_t}. We define \fairconst{h_t}{t} and
  \fairconstpr{h_t}{t} to be identical cost functions if
  the equal-cost constraints that they impose are identical, meaning one
  constraint is satisfied if and only if the other is satisfied. If this is not
  the case, then \fairconst{h_t}{t} and
  \fairconstpr{h_t}{t} are \emph{distinct},
  meaning that
  the equal-cost constraints are linearly independent for $\mu_1 \ne \mu_2$.
  Moreover, these are also linearly independent from the calibration constraints
  because by assumption, they both have nonzero coefficients for at least one of
  $(\fp{h_1}, \fn{h_1})$ and $(\fp{h_2}, \fn{h_2})$. As a result, we have four linearly
  independent constraints (2 from calibration and at least 2 equal-cost
  constraints) on 4 variables (\fp{h_1}, \fn{h_1}, \fp{h_2}, \fn{h_2}), meaning that these
  constraints yield a unique solution. From above, we know that all the
  constraints
  are simultaneously satisfied when $\fp{h_t} = \fn{h_t} = 0$ for $t = 1,2$, meaning
  that the perfect classifier is the only classifier for which they are
  simultaneously satisfied.
\end{proof}

\subsection{Approximate Impossibility Theorem}

Now, we will show that this impossibility result holds in an approximate sense --- i.e.
approximately satisfying the calibration and equal-cost constraints is only possible
if the classifiers approximately perfect.

Since the calibration and equal-cost constraints are all linear, let $A$ be the
matrix that encodes them. With two equal-cost constraints
\fairconstname{} and
$\fairconstname'$,
\[
  A = \begin{bmatrix}
    1 &-\frac{\mu_1}{1-\mu_1} & 0 & 0 \\
    0 & 0 & 1 &-\frac{\mu_2}{1-\mu_2} \\
    a_1 & b_1 & -a_2 & -b_2 \\
    a_1' & b_1' & -a_2' & -b_2'
  \end{bmatrix}.
\]
Note that the first two rows of $A$ encode the calibration conditions --- see \eqref{eqn:linear-calib-sup}.
The bottom two rows encode two equal-cost constraints.
Furthermore, let
\[
  \vec{q} = \left[\fp{h_1} \; \fn{h_1} \; \fp{h_2} \; \fn{h_2}\right]^\top.
\]
If all constraints are required to hold exactly, then we have $A \vec{q} = 0$.
Consider the case where the calibration and equal-cost constraints hold
approximately.
\begin{thm}[Generalized approximate impossibility result] \label{thm:approx-impossibility}
  Let $h_1$ and $h_2$ be classifiers with calibration $\delta_{cal}$ and
  cost difference at most $\delta_{cost}$ with respect to distinct cost functions
  $\fairconstname$ and $\fairconstname'$. Furthermore, assume that every entry
  of $A$ is rational with some common denominator $D$ and is upper bounded by
  some maximum value $M$. Then, there is a constant $L$ that depends on $D$ and
  $M$ such that
  \[
    \fp{h_t} \le L \cdot \max\left\{\frac{2\delta_{cal}}{1-\mu_1},
    \frac{2\delta_{cal}}{1-\mu_2}, \delta_{cost}\right\}
  \]
  and
  \[
    \fn{h_t} \le L \cdot \max\left\{\frac{2\delta_{cal}}{1-\mu_1},
    \frac{2\delta_{cal}}{1-\mu_2}, \delta_{cost}\right\}
  \]
  for $t = 1,2$.
  \label{thm:approx}
\end{thm}

\newcommand{\wA}{\widehat{A}}
\begin{proof}
  By Lemma \ref{lem:approx-linear-cal},
  \[
    \bigl|\mu_t \fn{h_t} - \left(1 - \mu_t \right) \fp{h_t} \bigr| \le 2
      \delta_{cal}.
  \]
  Since the first two rows in $A$ correspond to the calibration constraints, and
  the second to correspond to the equal-cost constraints, it must be the case
  that
  \[
    |A \vec{q}| \le \begin{bmatrix}
      \frac{2\delta_{cal}}{1-\mu_1} \\ \frac{2 \delta_{cal}}{1-\mu_2} \\
      \delta_{cost} \\ \delta_{cost}
    \end{bmatrix},
  \]
  i.e. the absolute value of each entry in $A\vec{q}$ is bounded by the vector
  on the right hand side. Let $\vec{\nu} = [ 2\delta_{cal}/(1-\mu_1) \; 2
  \delta_{cal}/(1-\mu_2) \; \delta_{cost} \; \delta_{cost}]^\top$.
  Let $\vec{s} = sign(A \vec{q})$, and multiply the
  $i$th row of $A$ by the $i$th entry of $\vec{s}$ to produce $\wA$
  This allows us to drop the absolute value, meaning we have
  \[
    \wA \vec{q} \le \vec{\nu}
  \]
  Furthermore, since $\fairconstname$ and $\fairconstname'$ were assumed to be
  distinct, $\wA$ is invertible, so this is equivalent to
  \[
    \vec{q} \le \wA^{-1}\vec{\nu}.
  \]
  Taking $\ell_{\infty}$ norms of both sides,
  \[
    \|\vec{q}\|_\infty \le \|\wA^{-1} \vec{\nu}\|_\infty \le \|
    \wA^{-1}\|_\infty \|\vec{\nu}\|_\infty.
  \]
  The $(i,j)$ entry of $\wA^{-1}$ can be expresed as $\wA_{ji}/\det(\wA)$,
  where $\wA_{ji}$ is the $(j,i)$ cofactor. Note that $\wA_{ji}$ is a $3 \times 3$
  determinant, so it is the sum of $6$ cubic polynomials in entries of $\wA$.
  However, since every $3 \times 3$ submatrix of $\wA$ has at least one $0$
  entry, only $4$ of those cubics can be nonnegative. By assumption, the maximum
  value of any entry of $\wA$ is $M$, so $|\wA_{ji}| \le 4 M^3$.

  We can lower bound $\det(\wA)$ by noting that since $\wA$ is not singular, its
  determinant is nonzero. However, because the determinant can be expressed as a
  $4 \times 4$ polynomial, and each term has common denominator $D$ by
  assumption, $|\det(\wA)| \ge 1/D^4$. As a result, $|\wA_{ji}/\det(\wA)| \le 4
  M^3 D^4$.

  Let $d_{ij}$ be the $(i,j)$ entry of $\wA$. We know that
  \[
    \|\wA^{-1}\|_\infty \le \max_j \sum_{i=1}^4 |d_{ij}| \le 16M^3 D^4 = L
  \]
  As a result,
  \[
    \|\vec{q}\|_\infty \le L \|\nu\|_\infty
  \]
  which proves the claim.
\end{proof}
%
Note that \autoref{thm:approx} is not intended to be a tight bound. It simply
shows that impossibility result degrades smoothly for approximate constraints.

